% REVISED sections for the literature review.
% Sections 1-3 (BSDE theory, numerical methods, MFG/MV-FBSDE) remain as-is.
% This replaces Section 4 (Empirical) and Section 5 (Conclusion).

% ============================================================
\section{Empirical Results and Financial Applications}\label{sec:lit-empirical-revised}

The translation of deep learning methods for BSDEs from mathematical
benchmarks to financial market applications is an active frontier,
with recent progress in three directions: deep BSDE solvers for
jump-diffusion processes, multi-agent game-theoretic models of dealer
markets, and the emergence of tacit collusion from decentralised
algorithmic learning.

\paragraph{Deep BSDE with jumps.}
The extension of the deep BSDE paradigm to jump-diffusion processes
is essential for financial applications where discrete events (order
executions, defaults) punctuate continuous price dynamics.
\citet{wang2023} developed deep learning methods for high-dimensional
fully nonlinear PIDEs and coupled FBSDEs with jumps (FBSDEJs), where
the jump-diffusion process is driven by both Brownian motion and a
compensated Poisson random measure.  Their method parameterises the
jump coefficients $U$ alongside the diffusion coefficient $Z$ as
neural networks, achieving convergence in up to 100 dimensions.
This is directly relevant to the dealer market-making problem, where
inventory changes via discrete Poisson execution events.
\citet{lu2024} applied deep reinforcement learning (actor-critic)
to multi-agent relative investment games in a jump-diffusion market,
demonstrating that neural network solutions converge to derived
constant Nash equilibria.  Their computational framework provides a
template for combining deep learning with stochastic games under
jump dynamics.

\paragraph{The Cont--Xiong dealer market model.}
The most comprehensive treatment of multi-agent market-making as a
stochastic differential game is \citet{cont2024}.  They model $N$
competing dealers in a continuous-time dealer market, formulated as
a stochastic differential game of intensity control with partial
information.  Each dealer's execution probability depends on their
own quote and those of all competitors through a logistic softmax
function (their equation~6), creating strategic interaction.
Competition is modelled as a Nash equilibrium, characterised by a
system of coupled Hamilton--Jacobi--Bellman equations (their
equation~28), for which existence is established via Brouwer's fixed
point theorem.  Collusion is modelled as a Pareto optimum of the
joint objective.  A key contribution is the fictitious play algorithm
(their Algorithm~1) for computing Nash equilibria by iterating
between policy evaluation (solving linear Bellman systems) and best
response calculation.

Critically, \citet{cont2024} show via multi-agent deep reinforcement
learning (Decentralised MADDPG) that competing market-making
algorithms, learning independently from their own trade outcomes
without any shared information, converge to spread levels strictly
above the Nash equilibrium --- a phenomenon they term ``tacit
collusion.''  This finding has significant regulatory implications:
the decentralised interaction of autonomous pricing algorithms may
produce outcomes resembling explicit cartel behaviour, even when no
coordination is intended.

\paragraph{Algorithmic collusion: broader evidence.}
The tacit collusion phenomenon extends beyond dealer markets.
\citet{calvano2020} first demonstrated in a Bertrand oligopoly that
Q-learning agents converge to supracompetitive prices, establishing
the ``algorithmic collusion'' paradigm.  Subsequent work has explored
the robustness and mechanisms of this phenomenon.
\citet{deng2025} conducted a comprehensive study of competitive and
collusive behaviours across multiple Bertrand oligopoly variants,
showing that the degree of algorithmic collusion depends critically
on the choice of RL algorithm: some algorithms (DQN, PPO) exhibit
stronger collusive tendencies than others.  This raises the question
of whether the collusion result is an artefact of the learning
algorithm or a genuine equilibrium phenomenon.

In the dealer market setting specifically, \citet{wang2025} study
multi-agent reinforcement learning for market making under
competition, showing that adaptive incentive control can support
competitive coexistence and mitigate collusive outcomes.  This
suggests that market design interventions, not just regulation, can
address algorithmic collusion.

\paragraph{The BSDE connection.}
A largely unexplored avenue is the application of deep BSDE methods
--- rather than model-free reinforcement learning --- to the
Cont--Xiong game.  The dealer's value function satisfies a backward
stochastic differential equation with Poisson jumps (BSDEJ), where
the jump coefficients encode the value changes on order execution.
The Cont--Xiong HJB system (their equation~28) is precisely the
stationary limit of this BSDEJ.  This observation connects the
game-theoretic market-making literature to the BSDE numerical methods
surveyed in Section~\ref{sec:lit-numerical}, creating the possibility
of importing convergence guarantees and architectural insights from
the deep BSDE literature into the multi-agent dealer market setting.

\paragraph{Recent surveys.}
\citet{hu2024survey} provide a comprehensive survey of recent
developments in machine learning methods for stochastic control and
games, covering deep BSDE methods, policy gradient approaches, and
mean-field game algorithms.  Their treatment highlights the growing
convergence between the BSDE, MFG, and multi-agent RL communities ---
precisely the intersection that this thesis occupies.

\paragraph{Synthesis.}
The empirical landscape reveals a specific gap: the Cont--Xiong
dealer market model provides a well-characterised benchmark with
exact Nash equilibria (Algorithm~1), exact Pareto optima
(Algorithm~0), and documented tacit collusion via MADDPG, yet no
work has applied deep BSDE methods to this model.  The deep BSDE
literature has demonstrated solvers for jump-diffusion processes in
high dimensions, and the game-theoretic literature has characterised
the equilibrium structure of competing market makers, but the two
have not been combined.  This thesis addresses this gap.


% ============================================================
\section{Conclusion and Open Problems}\label{sec:lit-conclusion-revised}

The literature reviewed here establishes four facts:
(i)~the mathematical theory of BSDEs with jumps and the
nonlinear Feynman--Kac theorem provide the rigorous foundation for
stochastic control problems with discrete events;
(ii)~deep learning methods can solve BSDEs and FBSDEs in high
dimensions without the curse of dimensionality;
(iii)~the Cont--Xiong dealer market model provides a precisely
characterised multi-agent game with known Nash equilibria, Pareto
optima, and documented tacit collusion;
but (iv)~no existing work applies deep BSDE methods to the
Cont--Xiong model or validates neural solvers against its exact
equilibrium benchmarks.

\paragraph{Scope of the present thesis.}
This thesis addresses the identified gap through four computational
contributions.

First, the Cont--Xiong dealer market model is formulated as a
backward stochastic differential equation with Poisson jumps
(BSDEJ), establishing that the coupled HJB system of
\citet{cont2024} is the stationary limit of a well-posed BSDEJ.
This connection places the numerical methods within the theoretical
framework of BSDE theory.

Second, a neural Bellman solver is developed that learns the value
function $V(q)$ on the discrete inventory grid and derives optimal
quotes from the first-order conditions of the BSDEJ generator.
The solver is validated against the exact Algorithm~1 of
\citet{cont2024}, achieving sub-1\% error at the Nash equilibrium
for $N = 1, 2, 5, 10$ competing dealers.  A neural fictitious play
algorithm is shown to converge to the Nash equilibrium in 10
iterations.

Third, the Pareto optimum (explicit collusion benchmark) is
computed via policy iteration on the joint state space, establishing
the upper bound on spreads.  The ordering Nash $<$ Monopolist $<$
Pareto is verified, providing reference bounds for the tacit
collusion analysis.

Fourth, tacit collusion is investigated using a decentralised
multi-agent deep deterministic policy gradient (MADDPG) algorithm
following \citet{cont2024}.  Multiple independent training rounds
are conducted to establish statistical significance of the finding
that learned spreads exceed the Nash equilibrium level.  The
information structure is varied from full information (converges to
Nash) to no information (converges to monopolist level), with
decentralised learning producing spreads between these bounds.

\paragraph{Methodological contribution.}
The primary methodological contribution is the validation of a
neural solver against the exact Cont--Xiong equilibrium.  Unlike
previous deep BSDE applications to financial models, where ground
truth is typically unavailable or limited to simple analytical
solutions, the Cont--Xiong Algorithm~1 provides an exact benchmark
at any number of dealers.  This enables rigorous error quantification
and architectural diagnosis --- determining not just \emph{whether}
the solver converges, but \emph{how accurately} and \emph{why}.

\paragraph{Open problems.}
Several directions remain for future work.  The extension to
common noise (systematic market shocks affecting all agents) would
require the deep signature approach of \citet{hu2025} or the
conditional distribution framework of \citet{carmona2018a}.  The
continuous-inventory formulation, where $q \in \mathbb{R}$ rather
than a discrete grid, would benefit from the full deep BSDE
architecture with learned $Z$-processes.  The multi-asset
generalisation, where dealers make markets in multiple correlated
instruments, would test the curse-of-dimensionality advantage of
neural methods over grid-based algorithms.  Finally, calibration
to real dealer market data would validate the economic content of
the model beyond its theoretical properties.
